\lysh{25326}{2}
\begin{alist}
\item Το πεδίο ορισμού της συνάρτησης $f$ είναι το σύνολο που περιέχει τις τιμές για τις οποίες ορίζεται η ρίζα, δηλαδή $x \ge 0$. Συνεπώς $D_f = [0, +\infty)$.

\item Βρίσκουμε την πρώτη παράγωγο της συνάρτησης.

Παραγωγίζουμε τη συνάρτηση ως άθροισμα δύο συναρτήσεων:

$f'(x) = (2x)' + (\sqrt{x})' = 2(x)' + (\sqrt{x})'$.

Άρα, $f'(x) = 2 \cdot 1 + \frac{1}{2\sqrt{x}} = 2 + \frac{1}{2\sqrt{x}}$, για κάθε $x \in (0, +\infty)$.

\item Από το β) ερώτημα έχουμε $f'(x) = 2 + \frac{1}{2\sqrt{x}}$.

Στην παραπάνω συνάρτηση $f'(x)$ βάζουμε όπου $x$ το $4$.

Άρα έχουμε $f'(4) = 2 + \frac{1}{2\sqrt{4}} = 2 + \frac{1}{4} = \frac{9}{4}$.
\end{alist}
\vspace{6mm}
\noindent\rule{\textwidth}{0.4pt}
\vspace{6mm}

